% !TeX TXS-program:compile = txs:///arara
% arara: pdflatex: {shell: no, synctex: no, interaction: batchmode}
% arara: pdflatex: {shell: no, synctex: no, interaction: batchmode}

\documentclass{article}
\usepackage[margin=1in]{geometry}
\usepackage{bera}
\usepackage{inconsolata}
\usepackage[T1]{fontenc}
\usepackage[scale=0.875]{cabin}
\usepackage{skill-icons}
\usepackage{xcolor}
\pagecolor{yellow!5}
\usepackage{longtable}
\usepackage{booktabs}
\usepackage{hyperref}
\usepackage{fancyvrb}
\usepackage{shortvrb}
\MakeShortVerb{\|}
\setlength{\parindent}{0pt}

\newenvironment{skilliconcase}
{%
  \footnotesize
  \begin{longtable}{cc@{\hspace{2em}}llr}%dark + light + name + macro + page
    \cmidrule[\heavyrulewidth]{1-5}% \toprule
    \bfseries Dark & \bfseries Light & \bfseries Name & \bfseries Direct command & \bfseries PDF page \\
    \cmidrule{1-5}
    \endhead
}%
{%
    \cmidrule[\heavyrulewidth]{1-5}
  \end{longtable}%
}

\NewDocumentCommand{\showcaseskillicons}{mm}{%name + page
  \skillicons[theme=dark,height=2em]{#1} & \skillicons[theme=light,height=2em]{#1} & \itshape #1 & \ttfamily \textbackslash skillicons\{#1\} & \ttfamily #2 \\
}

\begin{document}

\title{The \textsf{skill-icons} package (0.1a) -- 29/10/2026}
\author{%
  Cédric Pierquet\\%
  \url{https://github.com/LelouchFR/skill-icons} \\%
  \url{https://github.com/cpierquet/latex-packages/tree/main/skill-icons} (\LaTeX{} package)
}
\date{cpierquet -- at -- outlook . fr}\maketitle

This package provides \LaTeX{} support for the \textsf{Skill Icons} project (\url{https://github.com/LelouchFR/skill-icons}), a set of icons of programming languages, tools, frameworks and software, with a \textit{dark} and a \textit{light} version.

\medskip

\hspace*{5mm}\fbox{\begin{minipage}{0.75\linewidth}{\emph{A set of 837 icons (logos of languages, tools, frameworks\ldots), each one with a \textcolor{violet}{dark} and a \textcolor{violet}{light} version.}}\end{minipage}}

\medskip

To use \textsf{skill-icons} in your document, include the package with |\usepackage{skill-icons}|.

\medskip

An icon can be accessed using the icon name. A list of all included icons with their respective \textit{page} can be found at the end of this document.

\medskip

The individual icons are provided as pages of two PDF files (one for each theme) and are ultimately loaded using the |\includegraphics| command from the |graphicx| package.

\medskip

The icons are \textbf{colored} (logos with their own background): the current color of the text has no effect on them.

\section{Example}

\begin{Verbatim}[frame=single]
...
\usepackage{skill-icons}
...
\begin{document}
...
Inline \skillicons{python} version, and light \skillicons[theme=light]{python} version
...
\end{document}
\end{Verbatim}

Result: Inline \skillicons{python} version, and light \skillicons[theme=light]{python} version

\pagebreak

\section{Usage}

\subsection{Generic macro}

\begin{Verbatim}[frame=single]
\skillicons[height=...,(d)strut=...,theme=...]<includegraphics options>{name}
\end{Verbatim}

Key \texttt{[height=...]} can be:

\begin{itemize}
  \item \texttt{dauto} (by default): inline insertion with height/depth;
  \item \texttt{auto}: inline insertion with height/nodepth;
  \item \texttt{a length}: normal insertion with fixed height;
  \item \texttt{manual height/depth}: normal insertion with fixed height/depth.
\end{itemize}

The keys \texttt{[(d)strut=...]} can be used to set the \textit{reference} box for dimension calculations.

\medskip

The key \texttt{[theme=...]} can be \texttt{dark} (by default) or \texttt{light}, see the next subsection.

\medskip

The given height is the height of the \textit{full} box of the icon (the icon may include a small border).

\medskip

Example: Inline{\setlength\fboxsep{0pt}\fbox{\skillicons{python}}}version or inline {\setlength\fboxsep{0pt}\fbox{\skillicons{python}}} version

\pagebreak

\subsection{Dark and light themes}

Each icon exists in two versions, stored in two PDF files with the same page numbers: \texttt{skill-icons-dark-all} and \texttt{skill-icons-light-all}.

\medskip

\begin{itemize}
  \item the \textbf{dark} theme (by default): icons with a \textit{dark} background (or designed for it);
  \item the \textbf{light} theme: the same icons, with a \textit{light} background (or designed for it);
  \item some icons are \textit{neutral}: the original project provides only one version, which is then the same in both PDF files.
\end{itemize}

The theme can be chosen:

\begin{itemize}
  \item for an icon, with the key \texttt{[theme=dark]} or \texttt{[theme=light]};
  \item for the whole document, with the package option, \verb|\usepackage[light]{skill-icons}|, or at any time with \verb|\skilliconssetup{theme=light}| (the setting is global, and can be changed by another call).
\end{itemize}

\begin{Verbatim}[frame=single]
\skillicons[theme=dark]{github} \skillicons[theme=light]{github}
\skilliconssetup{theme=light}
\skillicons{github} \skillicons[theme=dark]{github}
\end{Verbatim}

{\Huge \skillicons[theme=dark]{github} \skillicons[theme=light]{github}
\skilliconssetup{theme=light}
\skillicons{github} \skillicons[theme=dark]{github}}
\skilliconssetup{theme=dark}

\pagebreak

\subsection{Examples}

\begin{Verbatim}[frame=single]
%normal version: inline with automatic height/depth (from qX characters)
\skillicons[height=dauto]{python} %or \skillicons{python}
\end{Verbatim}

{\Huge q\skillicons{python}X}

\begin{Verbatim}[frame=single]
%normal version: inline with automatic height/depth (from specified characters)
(q\skillicons[height=dauto,dstrut=(qB)]{python}B)
\end{Verbatim}

{\Huge (q\skillicons[height=dauto,dstrut=(qB)]{python}B)}

\begin{Verbatim}[frame=single]
%normal version: inline with automatic height (from X character)
\skillicons[height=auto]{python}
\end{Verbatim}

{\Huge X\skillicons[height=auto]{python}X}

\begin{Verbatim}[frame=single]
%normal version: inline with manual height/depth
\skillicons[height={0.8em/-0.4em}]{python}
\end{Verbatim}

{\Huge qb\skillicons[height={0.8em/-0.4em}]{python}Xz}

\begin{Verbatim}[frame=single]
%fixed height
\skillicons[height=2in]{python}
\end{Verbatim}

\skillicons[height=1.5in]{python}

\subsection{Used packages}

\texttt{graphicx} is loaded and used by the package (the key handling relies on \texttt{expl3}, which is part of the \LaTeX{} kernel since October 2020).

\subsection{License and attribution}

The package is released under the \LaTeX{} Project Public License v1.3c or later.

\medskip

The icons come from the \textsf{Skill Icons} project, released under the \textsf{MIT} license: the original project is \textsf{tandpfun/skill-icons} (\url{https://github.com/tandpfun/skill-icons}, \textcopyright{} 2022 tandpfun), and the icon set used here (with the dark and light versions) is the one of \textsf{LelouchFR/skill-icons} (\url{https://github.com/LelouchFR/skill-icons}, \textcopyright{} 2024 Baptiste Zahnow). The MIT license requires to keep the copyright notice and the permission notice in all copies of the icons.

\medskip

The icons are the logos of programming languages, tools, companies and software, which remain the \textbf{property of their respective owners}; they are provided to identify these products and their use may be subject to the guidelines of their owners.

\subsection{Naming}

The icon names are those of the original project, in lowercase and without separators (\texttt{python}, \texttt{githubactions}, \texttt{vscode}, \texttt{reactnative}\ldots). A package named \textsf{skillicons} also exists on CTAN: \textsf{skill-icons} is a distinct package, and all its files are prefixed by \texttt{skill-icons}.

\subsection{Bugs}

For bug reports and feature requests, report on the GitHub repository \url{https://github.com/cpierquet/latex-packages/issues}.

\newpage

\section{List of icons}

The \textsf{PDF} files are \texttt{skill-icons-dark-all} and \texttt{skill-icons-light-all} (same page numbers).

\begin{skilliconcase}
\showcaseskillicons{aave}{1}
\showcaseskillicons{ableton}{2}
\showcaseskillicons{acrobat}{3}
\showcaseskillicons{activitypub}{4}
\showcaseskillicons{actix}{5}
\showcaseskillicons{adobespark}{6}
\showcaseskillicons{adonis}{7}
\showcaseskillicons{aero}{8}
\showcaseskillicons{affinity}{9}
\showcaseskillicons{aftereffects}{10}
\showcaseskillicons{agno}{11}
\showcaseskillicons{aiogram}{12}
\showcaseskillicons{airbyte}{13}
\showcaseskillicons{airflow}{14}
\showcaseskillicons{aiscript}{15}
\showcaseskillicons{alacritty}{16}
\showcaseskillicons{alchemy}{17}
\showcaseskillicons{alpinejs}{18}
\showcaseskillicons{amplify}{19}
\showcaseskillicons{anaconda}{20}
\showcaseskillicons{android}{21}
\showcaseskillicons{androidstudio}{22}
\showcaseskillicons{angular}{23}
\showcaseskillicons{animate}{24}
\showcaseskillicons{animejs}{25}
\showcaseskillicons{anki}{26}
\showcaseskillicons{ansible}{27}
\showcaseskillicons{ansys}{28}
\showcaseskillicons{antdesign}{29}
\showcaseskillicons{antigravity}{30}
\showcaseskillicons{apache}{31}
\showcaseskillicons{apeworx}{32}
\showcaseskillicons{apexcharts}{33}
\showcaseskillicons{api}{34}
\showcaseskillicons{apidog}{35}
\showcaseskillicons{apigateway}{36}
\showcaseskillicons{apollo}{37}
\showcaseskillicons{appactive}{38}
\showcaseskillicons{appcode}{39}
\showcaseskillicons{apple}{40}
\showcaseskillicons{appstore}{41}
\showcaseskillicons{apptainer}{42}
\showcaseskillicons{appwrite}{43}
\showcaseskillicons{aqua}{44}
\showcaseskillicons{arcbrowser}{45}
\showcaseskillicons{arch}{46}
\showcaseskillicons{arcjet}{47}
\showcaseskillicons{arduino}{48}
\showcaseskillicons{argocd}{49}
\showcaseskillicons{arrow}{50}
\showcaseskillicons{aseprite}{51}
\showcaseskillicons{assembly}{52}
\showcaseskillicons{astro}{53}
\showcaseskillicons{athena}{54}
\showcaseskillicons{atom}{55}
\showcaseskillicons{audacity}{56}
\showcaseskillicons{audition}{57}
\showcaseskillicons{aurora}{58}
\showcaseskillicons{authenticator}{59}
\showcaseskillicons{authjs}{60}
\showcaseskillicons{autocad}{61}
\showcaseskillicons{avaloniaui}{62}
\showcaseskillicons{aws}{63}
\showcaseskillicons{axios}{64}
\showcaseskillicons{azul}{65}
\showcaseskillicons{azure}{66}
\showcaseskillicons{azuredevops}{67}
\showcaseskillicons{babel}{68}
\showcaseskillicons{balancer}{69}
\showcaseskillicons{barbajs}{70}
\showcaseskillicons{bash}{71}
\showcaseskillicons{beam}{72}
\showcaseskillicons{beeceptor}{73}
\showcaseskillicons{behance}{74}
\showcaseskillicons{behat}{75}
\showcaseskillicons{betterauth}{76}
\showcaseskillicons{bevy}{77}
\showcaseskillicons{bigquery}{78}
\showcaseskillicons{biome}{79}
\showcaseskillicons{bitbucket}{80}
\showcaseskillicons{bitrix24}{81}
\showcaseskillicons{blazor}{82}
\showcaseskillicons{blender}{83}
\showcaseskillicons{bluesky}{84}
\showcaseskillicons{bokeh}{85}
\showcaseskillicons{bootstrap}{86}
\showcaseskillicons{brave}{87}
\showcaseskillicons{breeze}{88}
\showcaseskillicons{bridge}{89}
\showcaseskillicons{bsd}{90}
\showcaseskillicons{btlo}{91}
\showcaseskillicons{bulma}{92}
\showcaseskillicons{bun}{93}
\showcaseskillicons{burn}{94}
\showcaseskillicons{burpsuite}{95}
\showcaseskillicons{c}{96}
\showcaseskillicons{cachyos}{97}
\showcaseskillicons{caddy}{98}
\showcaseskillicons{cairo}{99}
\showcaseskillicons{canva}{100}
\showcaseskillicons{capacitor}{101}
\showcaseskillicons{capture}{102}
\showcaseskillicons{cashier}{103}
\showcaseskillicons{cassandra}{104}
\showcaseskillicons{catppuccin}{105}
\showcaseskillicons{celerdata}{106}
\showcaseskillicons{celery}{107}
\showcaseskillicons{centos}{108}
\showcaseskillicons{chainlink}{109}
\showcaseskillicons{chakraui}{110}
\showcaseskillicons{chaosblade}{111}
\showcaseskillicons{characteranimator}{112}
\showcaseskillicons{chartjs}{113}
\showcaseskillicons{chatgpt}{114}
\showcaseskillicons{chi}{115}
\showcaseskillicons{chrome}{116}
\showcaseskillicons{chromedevtools}{117}
\showcaseskillicons{chromium}{118}
\showcaseskillicons{circleci}{119}
\showcaseskillicons{claude}{120}
\showcaseskillicons{clerk}{121}
\showcaseskillicons{clickhouse}{122}
\showcaseskillicons{clickup}{123}
\showcaseskillicons{clion}{124}
\showcaseskillicons{clojure}{125}
\showcaseskillicons{cloudflare}{126}
\showcaseskillicons{cloudformation}{127}
\showcaseskillicons{cloudfront}{128}
\showcaseskillicons{cloudwatch}{129}
\showcaseskillicons{cmake}{130}
\showcaseskillicons{cockroachdb}{131}
\showcaseskillicons{codeberg}{132}
\showcaseskillicons{codeblocks}{133}
\showcaseskillicons{codeigniter}{134}
\showcaseskillicons{codepen}{135}
\showcaseskillicons{coffeescript}{136}
\showcaseskillicons{commercetools}{137}
\showcaseskillicons{composer}{138}
\showcaseskillicons{confluence}{139}
\showcaseskillicons{confluent}{140}
\showcaseskillicons{consul}{141}
\showcaseskillicons{contentful}{142}
\showcaseskillicons{cpp}{143}
\showcaseskillicons{creativecloud}{144}
\showcaseskillicons{crewai}{145}
\showcaseskillicons{crystal}{146}
\showcaseskillicons{cs}{147}
\showcaseskillicons{css}{148}
\showcaseskillicons{cuda}{149}
\showcaseskillicons{cursor}{150}
\showcaseskillicons{cypress}{151}
\showcaseskillicons{d}{152}
\showcaseskillicons{d3}{153}
\showcaseskillicons{daft}{154}
\showcaseskillicons{dailydev}{155}
\showcaseskillicons{daisyui}{156}
\showcaseskillicons{dapper}{157}
\showcaseskillicons{dart}{158}
\showcaseskillicons{dask}{159}
\showcaseskillicons{databricks}{160}
\showcaseskillicons{datadog}{161}
\showcaseskillicons{datagrip}{162}
\showcaseskillicons{dataspell}{163}
\showcaseskillicons{davinci}{164}
\showcaseskillicons{dbeaver}{165}
\showcaseskillicons{dbtlabs}{166}
\showcaseskillicons{debian}{167}
\showcaseskillicons{deepseek}{168}
\showcaseskillicons{defold}{169}
\showcaseskillicons{delta}{170}
\showcaseskillicons{deltars}{171}
\showcaseskillicons{deno}{172}
\showcaseskillicons{desmos}{173}
\showcaseskillicons{devto}{174}
\showcaseskillicons{digitalocean}{175}
\showcaseskillicons{dimension}{176}
\showcaseskillicons{directus}{177}
\showcaseskillicons{discord}{178}
\showcaseskillicons{discordbots}{179}
\showcaseskillicons{discordjs}{180}
\showcaseskillicons{django}{181}
\showcaseskillicons{djangorestframework}{182}
\showcaseskillicons{dn42}{183}
\showcaseskillicons{docker}{184}
\showcaseskillicons{docksal}{185}
\showcaseskillicons{docsify}{186}
\showcaseskillicons{doctrine}{187}
\showcaseskillicons{doris}{188}
\showcaseskillicons{dotnet}{189}
\showcaseskillicons{dreamweaver}{190}
\showcaseskillicons{dremio}{191}
\showcaseskillicons{drizzle}{192}
\showcaseskillicons{drupal}{193}
\showcaseskillicons{dubbo}{194}
\showcaseskillicons{duckdb}{195}
\showcaseskillicons{duckduckgo}{196}
\showcaseskillicons{dusk}{197}
\showcaseskillicons{dynamodb}{198}
\showcaseskillicons{ec2}{199}
\showcaseskillicons{echo}{200}
\showcaseskillicons{eclipse}{201}
\showcaseskillicons{ecr}{202}
\showcaseskillicons{ecs}{203}
\showcaseskillicons{edge}{204}
\showcaseskillicons{ejs}{205}
\showcaseskillicons{eks}{206}
\showcaseskillicons{elasticbeanstalk}{207}
\showcaseskillicons{elasticsearch}{208}
\showcaseskillicons{elb}{209}
\showcaseskillicons{electron}{210}
\showcaseskillicons{element}{211}
\showcaseskillicons{elementor}{212}
\showcaseskillicons{elementplus}{213}
\showcaseskillicons{elixir}{214}
\showcaseskillicons{elm}{215}
\showcaseskillicons{elysia}{216}
\showcaseskillicons{emacs}{217}
\showcaseskillicons{ember}{218}
\showcaseskillicons{emotion}{219}
\showcaseskillicons{emr}{220}
\showcaseskillicons{envoyer}{221}
\showcaseskillicons{erlang}{222}
\showcaseskillicons{eslint}{223}
\showcaseskillicons{etcd}{224}
\showcaseskillicons{ethereum}{225}
\showcaseskillicons{eventbridge}{226}
\showcaseskillicons{excel}{227}
\showcaseskillicons{expo}{228}
\showcaseskillicons{expressjs}{229}
\showcaseskillicons{fabric}{230}
\showcaseskillicons{fabricmc}{231}
\showcaseskillicons{facebook}{232}
\showcaseskillicons{fargate}{233}
\showcaseskillicons{fastai}{234}
\showcaseskillicons{fastapi}{235}
\showcaseskillicons{fastlane}{236}
\showcaseskillicons{fediverse}{237}
\showcaseskillicons{fedora}{238}
\showcaseskillicons{ffmpeg}{239}
\showcaseskillicons{fiber}{240}
\showcaseskillicons{figma}{241}
\showcaseskillicons{filament}{242}
\showcaseskillicons{filmora}{243}
\showcaseskillicons{firebase}{244}
\showcaseskillicons{firefox}{245}
\showcaseskillicons{fiverr}{246}
\showcaseskillicons{fivetran}{247}
\showcaseskillicons{flameshot}{248}
\showcaseskillicons{flask}{249}
\showcaseskillicons{fleet}{250}
\showcaseskillicons{flink}{251}
\showcaseskillicons{flutter}{252}
\showcaseskillicons{flutterflow}{253}
\showcaseskillicons{flyio}{254}
\showcaseskillicons{fonts}{255}
\showcaseskillicons{forge}{256}
\showcaseskillicons{forgejo}{257}
\showcaseskillicons{forgemc}{258}
\showcaseskillicons{forth}{259}
\showcaseskillicons{fortran}{260}
\showcaseskillicons{foundry}{261}
\showcaseskillicons{framer}{262}
\showcaseskillicons{frankenphp}{263}
\showcaseskillicons{freecad}{264}
\showcaseskillicons{freecodecamp}{265}
\showcaseskillicons{freelancer}{266}
\showcaseskillicons{fresco}{267}
\showcaseskillicons{fresh}{268}
\showcaseskillicons{fuse}{269}
\showcaseskillicons{gamemakerstudio}{270}
\showcaseskillicons{ganache}{271}
\showcaseskillicons{gatsby}{272}
\showcaseskillicons{gcp}{273}
\showcaseskillicons{gdevelop}{274}
\showcaseskillicons{gemini}{275}
\showcaseskillicons{gentoo}{276}
\showcaseskillicons{gherkin}{277}
\showcaseskillicons{ghostty}{278}
\showcaseskillicons{gimp}{279}
\showcaseskillicons{gin}{280}
\showcaseskillicons{git}{281}
\showcaseskillicons{gitbash}{282}
\showcaseskillicons{gitea}{283}
\showcaseskillicons{github}{284}
\showcaseskillicons{githubactions}{285}
\showcaseskillicons{githubcopilot}{286}
\showcaseskillicons{githubdesktop}{287}
\showcaseskillicons{githubpages}{288}
\showcaseskillicons{gitkraken}{289}
\showcaseskillicons{gitlab}{290}
\showcaseskillicons{gleam}{291}
\showcaseskillicons{glue}{292}
\showcaseskillicons{gmail}{293}
\showcaseskillicons{gmx}{294}
\showcaseskillicons{gnome}{295}
\showcaseskillicons{godot}{296}
\showcaseskillicons{goland}{297}
\showcaseskillicons{golang}{298}
\showcaseskillicons{googleanalytics}{299}
\showcaseskillicons{googleappsscript}{300}
\showcaseskillicons{googlecolab}{301}
\showcaseskillicons{googleplay}{302}
\showcaseskillicons{googleplayconsole}{303}
\showcaseskillicons{gorm}{304}
\showcaseskillicons{gradio}{305}
\showcaseskillicons{gradle}{306}
\showcaseskillicons{grafana}{307}
\showcaseskillicons{grails}{308}
\showcaseskillicons{granica}{309}
\showcaseskillicons{graphite}{310}
\showcaseskillicons{graphql}{311}
\showcaseskillicons{grok}{312}
\showcaseskillicons{gromacs}{313}
\showcaseskillicons{groq}{314}
\showcaseskillicons{grpc}{315}
\showcaseskillicons{grunt}{316}
\showcaseskillicons{gsap}{317}
\showcaseskillicons{gtk}{318}
\showcaseskillicons{gulp}{319}
\showcaseskillicons{hackerrank}{320}
\showcaseskillicons{hackthebox}{321}
\showcaseskillicons{hadoop}{322}
\showcaseskillicons{hardhat}{323}
\showcaseskillicons{haskell}{324}
\showcaseskillicons{haxe}{325}
\showcaseskillicons{haxeflixel}{326}
\showcaseskillicons{helia}{327}
\showcaseskillicons{helix}{328}
\showcaseskillicons{helm}{329}
\showcaseskillicons{herd}{330}
\showcaseskillicons{heroku}{331}
\showcaseskillicons{hexo}{332}
\showcaseskillicons{hibernate}{333}
\showcaseskillicons{higress}{334}
\showcaseskillicons{hive}{335}
\showcaseskillicons{holyc}{336}
\showcaseskillicons{hono}{337}
\showcaseskillicons{horizon}{338}
\showcaseskillicons{html}{339}
\showcaseskillicons{htmx}{340}
\showcaseskillicons{htop}{341}
\showcaseskillicons{hudi}{342}
\showcaseskillicons{huggingface}{343}
\showcaseskillicons{hugo}{344}
\showcaseskillicons{hydrogen}{345}
\showcaseskillicons{hyprland}{346}
\showcaseskillicons{i3}{347}
\showcaseskillicons{iceberg}{348}
\showcaseskillicons{iced}{349}
\showcaseskillicons{idea}{350}
\showcaseskillicons{ignite}{351}
\showcaseskillicons{illustrator}{352}
\showcaseskillicons{immuta}{353}
\showcaseskillicons{impala}{354}
\showcaseskillicons{incopy}{355}
\showcaseskillicons{indesign}{356}
\showcaseskillicons{inertia}{357}
\showcaseskillicons{informatica}{358}
\showcaseskillicons{infura}{359}
\showcaseskillicons{inkscape}{360}
\showcaseskillicons{insomnia}{361}
\showcaseskillicons{instagram}{362}
\showcaseskillicons{integrations}{363}
\showcaseskillicons{ipfs}{364}
\showcaseskillicons{itchio}{365}
\showcaseskillicons{jaeger}{366}
\showcaseskillicons{jamovi}{367}
\showcaseskillicons{java}{368}
\showcaseskillicons{javascript}{369}
\showcaseskillicons{jax}{370}
\showcaseskillicons{jekyll}{371}
\showcaseskillicons{jenkins}{372}
\showcaseskillicons{jest}{373}
\showcaseskillicons{jetpackcompose}{374}
\showcaseskillicons{jetstream}{375}
\showcaseskillicons{jira}{376}
\showcaseskillicons{joomla}{377}
\showcaseskillicons{jqlang}{378}
\showcaseskillicons{jquery}{379}
\showcaseskillicons{json}{380}
\showcaseskillicons{julia}{381}
\showcaseskillicons{junit}{382}
\showcaseskillicons{jupyter}{383}
\showcaseskillicons{jwt}{384}
\showcaseskillicons{k3s}{385}
\showcaseskillicons{kafka}{386}
\showcaseskillicons{kaggle}{387}
\showcaseskillicons{kakoune}{388}
\showcaseskillicons{kali}{389}
\showcaseskillicons{karma}{390}
\showcaseskillicons{kde}{391}
\showcaseskillicons{kdenlive}{392}
\showcaseskillicons{keycloak}{393}
\showcaseskillicons{keydb}{394}
\showcaseskillicons{kibana}{395}
\showcaseskillicons{kitty}{396}
\showcaseskillicons{kong}{397}
\showcaseskillicons{kotlin}{398}
\showcaseskillicons{ktor}{399}
\showcaseskillicons{kubernetes}{400}
\showcaseskillicons{kubevela}{401}
\showcaseskillicons{lambda}{402}
\showcaseskillicons{lancedb}{403}
\showcaseskillicons{lando}{404}
\showcaseskillicons{langchain}{405}
\showcaseskillicons{laravel}{406}
\showcaseskillicons{laravelspark}{407}
\showcaseskillicons{latex}{408}
\showcaseskillicons{lazyvim}{409}
\showcaseskillicons{leaflet}{410}
\showcaseskillicons{leetcode}{411}
\showcaseskillicons{less}{412}
\showcaseskillicons{libreoffice}{413}
\showcaseskillicons{librewolf}{414}
\showcaseskillicons{libsql}{415}
\showcaseskillicons{lighthouse}{416}
\showcaseskillicons{lightning}{417}
\showcaseskillicons{lightroom}{418}
\showcaseskillicons{lightroomclassic}{419}
\showcaseskillicons{linkedin}{420}
\showcaseskillicons{linux}{421}
\showcaseskillicons{liquidsoap}{422}
\showcaseskillicons{lit}{423}
\showcaseskillicons{litestar}{424}
\showcaseskillicons{litmus}{425}
\showcaseskillicons{livewire}{426}
\showcaseskillicons{llamaindex}{427}
\showcaseskillicons{logto}{428}
\showcaseskillicons{looker}{429}
\showcaseskillicons{lottielab}{430}
\showcaseskillicons{lua}{431}
\showcaseskillicons{luau}{432}
\showcaseskillicons{lucidchart}{433}
\showcaseskillicons{lunacy}{434}
\showcaseskillicons{lxc}{435}
\showcaseskillicons{lynxjs}{436}
\showcaseskillicons{macos}{437}
\showcaseskillicons{manim}{438}
\showcaseskillicons{manjaro}{439}
\showcaseskillicons{mariadb}{440}
\showcaseskillicons{markdown}{441}
\showcaseskillicons{mastodon}{442}
\showcaseskillicons{materialui}{443}
\showcaseskillicons{matlab}{444}
\showcaseskillicons{matplotlib}{445}
\showcaseskillicons{maven}{446}
\showcaseskillicons{max8}{447}
\showcaseskillicons{max9}{448}
\showcaseskillicons{mcp}{449}
\showcaseskillicons{mdbook}{450}
\showcaseskillicons{mediaencoder}{451}
\showcaseskillicons{medium}{452}
\showcaseskillicons{mermaid}{453}
\showcaseskillicons{metabase}{454}
\showcaseskillicons{meteorjs}{455}
\showcaseskillicons{micropython}{456}
\showcaseskillicons{microsoftcopilot}{457}
\showcaseskillicons{miktex}{458}
\showcaseskillicons{millionjs}{459}
\showcaseskillicons{milvus}{460}
\showcaseskillicons{mindsdb}{461}
\showcaseskillicons{mint}{462}
\showcaseskillicons{miro}{463}
\showcaseskillicons{misskey}{464}
\showcaseskillicons{mistral}{465}
\showcaseskillicons{mjml}{466}
\showcaseskillicons{ml5}{467}
\showcaseskillicons{mlflow}{468}
\showcaseskillicons{mobx}{469}
\showcaseskillicons{mocha}{470}
\showcaseskillicons{modelviewer}{471}
\showcaseskillicons{mojo}{472}
\showcaseskillicons{mongodb}{473}
\showcaseskillicons{mongoose}{474}
\showcaseskillicons{mux}{475}
\showcaseskillicons{mysql}{476}
\showcaseskillicons{n8n}{477}
\showcaseskillicons{nacos}{478}
\showcaseskillicons{neo4j}{479}
\showcaseskillicons{neoforge}{480}
\showcaseskillicons{neon}{481}
\showcaseskillicons{neovim}{482}
\showcaseskillicons{nestjs}{483}
\showcaseskillicons{netlify}{484}
\showcaseskillicons{nextflow}{485}
\showcaseskillicons{nextjs}{486}
\showcaseskillicons{nginx}{487}
\showcaseskillicons{ngrok}{488}
\showcaseskillicons{ngrx}{489}
\showcaseskillicons{nim}{490}
\showcaseskillicons{nixos}{491}
\showcaseskillicons{nodejs}{492}
\showcaseskillicons{notepadpp}{493}
\showcaseskillicons{notion}{494}
\showcaseskillicons{nova}{495}
\showcaseskillicons{npm}{496}
\showcaseskillicons{numpy}{497}
\showcaseskillicons{nunjucks}{498}
\showcaseskillicons{nuxtjs}{499}
\showcaseskillicons{nvidia}{500}
\showcaseskillicons{obs}{501}
\showcaseskillicons{obsidian}{502}
\showcaseskillicons{ocaml}{503}
\showcaseskillicons{octane}{504}
\showcaseskillicons{octave}{505}
\showcaseskillicons{odin}{506}
\showcaseskillicons{ollama}{507}
\showcaseskillicons{onedrive}{508}
\showcaseskillicons{onehouse}{509}
\showcaseskillicons{onenote}{510}
\showcaseskillicons{opencode}{511}
\showcaseskillicons{opencv}{512}
\showcaseskillicons{openmm}{513}
\showcaseskillicons{opensergo}{514}
\showcaseskillicons{openshift}{515}
\showcaseskillicons{opensource}{516}
\showcaseskillicons{openstack}{517}
\showcaseskillicons{opentelemetry}{518}
\showcaseskillicons{openzeppelin}{519}
\showcaseskillicons{opera}{520}
\showcaseskillicons{oracle}{521}
\showcaseskillicons{orchid}{522}
\showcaseskillicons{outlook}{523}
\showcaseskillicons{overleaf}{524}
\showcaseskillicons{p4}{525}
\showcaseskillicons{p5js}{526}
\showcaseskillicons{pail}{527}
\showcaseskillicons{pancakeswap}{528}
\showcaseskillicons{pandas}{529}
\showcaseskillicons{papertrail}{530}
\showcaseskillicons{payload}{531}
\showcaseskillicons{pbi}{532}
\showcaseskillicons{pennant}{533}
\showcaseskillicons{perl}{534}
\showcaseskillicons{phaser}{535}
\showcaseskillicons{photoshop}{536}
\showcaseskillicons{photoshopclassic}{537}
\showcaseskillicons{photoshopexpress}{538}
\showcaseskillicons{php}{539}
\showcaseskillicons{phpstan}{540}
\showcaseskillicons{phpstorm}{541}
\showcaseskillicons{picocss}{542}
\showcaseskillicons{pinecone}{543}
\showcaseskillicons{pinescript}{544}
\showcaseskillicons{pinia}{545}
\showcaseskillicons{pint}{546}
\showcaseskillicons{pkl}{547}
\showcaseskillicons{plan9}{548}
\showcaseskillicons{planetscale}{549}
\showcaseskillicons{platformio}{550}
\showcaseskillicons{playcanvas}{551}
\showcaseskillicons{playfab}{552}
\showcaseskillicons{playwright}{553}
\showcaseskillicons{plotly}{554}
\showcaseskillicons{plsql}{555}
\showcaseskillicons{pm2}{556}
\showcaseskillicons{pnpm}{557}
\showcaseskillicons{pocketbase}{558}
\showcaseskillicons{podman}{559}
\showcaseskillicons{polar}{560}
\showcaseskillicons{polars}{561}
\showcaseskillicons{popos}{562}
\showcaseskillicons{portainer}{563}
\showcaseskillicons{portfolio}{564}
\showcaseskillicons{postcss}{565}
\showcaseskillicons{postgresql}{566}
\showcaseskillicons{postman}{567}
\showcaseskillicons{powerautomate}{568}
\showcaseskillicons{powerpoint}{569}
\showcaseskillicons{powershell}{570}
\showcaseskillicons{powertoys}{571}
\showcaseskillicons{preact}{572}
\showcaseskillicons{prelude}{573}
\showcaseskillicons{premiere}{574}
\showcaseskillicons{premiererush}{575}
\showcaseskillicons{presto}{576}
\showcaseskillicons{prettier}{577}
\showcaseskillicons{primeng}{578}
\showcaseskillicons{primereact}{579}
\showcaseskillicons{primevue}{580}
\showcaseskillicons{prisma}{581}
\showcaseskillicons{processing}{582}
\showcaseskillicons{prometheus}{583}
\showcaseskillicons{prompts}{584}
\showcaseskillicons{proton}{585}
\showcaseskillicons{proxmox}{586}
\showcaseskillicons{pug}{587}
\showcaseskillicons{pulsar}{588}
\showcaseskillicons{pulse}{589}
\showcaseskillicons{pulumi}{590}
\showcaseskillicons{puppeteer}{591}
\showcaseskillicons{puppygraph}{592}
\showcaseskillicons{putty}{593}
\showcaseskillicons{pwa}{594}
\showcaseskillicons{pycharm}{595}
\showcaseskillicons{pydantic}{596}
\showcaseskillicons{pygame}{597}
\showcaseskillicons{pypi}{598}
\showcaseskillicons{pyspark}{599}
\showcaseskillicons{pytest}{600}
\showcaseskillicons{python}{601}
\showcaseskillicons{pytorch}{602}
\showcaseskillicons{pyxel}{603}
\showcaseskillicons{qdrant}{604}
\showcaseskillicons{qemu}{605}
\showcaseskillicons{qodana}{606}
\showcaseskillicons{qt}{607}
\showcaseskillicons{quarkus}{608}
\showcaseskillicons{qubesos}{609}
\showcaseskillicons{querydsl}{610}
\showcaseskillicons{quiltmc}{611}
\showcaseskillicons{r}{612}
\showcaseskillicons{rabbitmq}{613}
\showcaseskillicons{radix}{614}
\showcaseskillicons{rails}{615}
\showcaseskillicons{railway}{616}
\showcaseskillicons{rancher}{617}
\showcaseskillicons{raspberrypi}{618}
\showcaseskillicons{ratatui}{619}
\showcaseskillicons{ray}{620}
\showcaseskillicons{rclone}{621}
\showcaseskillicons{rds}{622}
\showcaseskillicons{react}{623}
\showcaseskillicons{reactbootstrap}{624}
\showcaseskillicons{reactivex}{625}
\showcaseskillicons{reactlynx}{626}
\showcaseskillicons{reactnative}{627}
\showcaseskillicons{reactos}{628}
\showcaseskillicons{reactquery}{629}
\showcaseskillicons{reactrouter}{630}
\showcaseskillicons{recoil}{631}
\showcaseskillicons{reddit}{632}
\showcaseskillicons{redhat}{633}
\showcaseskillicons{redis}{634}
\showcaseskillicons{redshift}{635}
\showcaseskillicons{redux}{636}
\showcaseskillicons{regex}{637}
\showcaseskillicons{remix}{638}
\showcaseskillicons{render}{639}
\showcaseskillicons{renpy}{640}
\showcaseskillicons{replit}{641}
\showcaseskillicons{resend}{642}
\showcaseskillicons{resharper}{643}
\showcaseskillicons{restructuredtext}{644}
\showcaseskillicons{reverb}{645}
\showcaseskillicons{revolt}{646}
\showcaseskillicons{rider}{647}
\showcaseskillicons{robloxstudio}{648}
\showcaseskillicons{rocket}{649}
\showcaseskillicons{rocketmq}{650}
\showcaseskillicons{rollupjs}{651}
\showcaseskillicons{ros}{652}
\showcaseskillicons{rubocop}{653}
\showcaseskillicons{ruby}{654}
\showcaseskillicons{rubymine}{655}
\showcaseskillicons{rust}{656}
\showcaseskillicons{rustrover}{657}
\showcaseskillicons{s3}{658}
\showcaseskillicons{safari}{659}
\showcaseskillicons{sail}{660}
\showcaseskillicons{salesforce}{661}
\showcaseskillicons{sanctum}{662}
\showcaseskillicons{sanity}{663}
\showcaseskillicons{sas}{664}
\showcaseskillicons{sass}{665}
\showcaseskillicons{scala}{666}
\showcaseskillicons{scikitlearn}{667}
\showcaseskillicons{scipy}{668}
\showcaseskillicons{scout}{669}
\showcaseskillicons{scratch}{670}
\showcaseskillicons{seaborn}{671}
\showcaseskillicons{seata}{672}
\showcaseskillicons{selenium}{673}
\showcaseskillicons{sentinel}{674}
\showcaseskillicons{sentry}{675}
\showcaseskillicons{sequelize}{676}
\showcaseskillicons{ses}{677}
\showcaseskillicons{session}{678}
\showcaseskillicons{shadcn}{679}
\showcaseskillicons{sharepoint}{680}
\showcaseskillicons{shopify}{681}
\showcaseskillicons{signal}{682}
\showcaseskillicons{skeletonui}{683}
\showcaseskillicons{sketchup}{684}
\showcaseskillicons{skywalking}{685}
\showcaseskillicons{slack}{686}
\showcaseskillicons{snowflake}{687}
\showcaseskillicons{snyk}{688}
\showcaseskillicons{socialite}{689}
\showcaseskillicons{socketio}{690}
\showcaseskillicons{solana}{691}
\showcaseskillicons{solidity}{692}
\showcaseskillicons{solidjs}{693}
\showcaseskillicons{sonarqube}{694}
\showcaseskillicons{spark}{695}
\showcaseskillicons{sparksql}{696}
\showcaseskillicons{sphinx}{697}
\showcaseskillicons{spring}{698}
\showcaseskillicons{springbatch}{699}
\showcaseskillicons{springdatajpa}{700}
\showcaseskillicons{springsecurity}{701}
\showcaseskillicons{sqlalchemy}{702}
\showcaseskillicons{sqlite}{703}
\showcaseskillicons{sqlserver}{704}
\showcaseskillicons{sqs}{705}
\showcaseskillicons{stackoverflow}{706}
\showcaseskillicons{stan}{707}
\showcaseskillicons{starburst}{708}
\showcaseskillicons{starrocks}{709}
\showcaseskillicons{steam}{710}
\showcaseskillicons{stock}{711}
\showcaseskillicons{storyblok}{712}
\showcaseskillicons{storybook}{713}
\showcaseskillicons{strapi}{714}
\showcaseskillicons{streamlit}{715}
\showcaseskillicons{stripe}{716}
\showcaseskillicons{styledcomponents}{717}
\showcaseskillicons{stylelint}{718}
\showcaseskillicons{stylus}{719}
\showcaseskillicons{sublime}{720}
\showcaseskillicons{supabase}{721}
\showcaseskillicons{surrealdb}{722}
\showcaseskillicons{sushiswap}{723}
\showcaseskillicons{svelte}{724}
\showcaseskillicons{svg}{725}
\showcaseskillicons{svn}{726}
\showcaseskillicons{swagger}{727}
\showcaseskillicons{swift}{728}
\showcaseskillicons{symfony}{729}
\showcaseskillicons{systemd}{730}
\showcaseskillicons{t3}{731}
\showcaseskillicons{tableau}{732}
\showcaseskillicons{taiga}{733}
\showcaseskillicons{tailscale}{734}
\showcaseskillicons{tailsos}{735}
\showcaseskillicons{tailwindcss}{736}
\showcaseskillicons{tallyprime}{737}
\showcaseskillicons{tanstack}{738}
\showcaseskillicons{tauri}{739}
\showcaseskillicons{teams}{740}
\showcaseskillicons{tecton}{741}
\showcaseskillicons{telegram}{742}
\showcaseskillicons{telescope}{743}
\showcaseskillicons{tensorflow}{744}
\showcaseskillicons{terminal}{745}
\showcaseskillicons{terraform}{746}
\showcaseskillicons{testinglibrary}{747}
\showcaseskillicons{texmaker}{748}
\showcaseskillicons{threejs}{749}
\showcaseskillicons{thunderbird}{750}
\showcaseskillicons{thunkable}{751}
\showcaseskillicons{tidb}{752}
\showcaseskillicons{tmux}{753}
\showcaseskillicons{tokiors}{754}
\showcaseskillicons{tomcat}{755}
\showcaseskillicons{toml}{756}
\showcaseskillicons{tor}{757}
\showcaseskillicons{touchdesigner}{758}
\showcaseskillicons{trino}{759}
\showcaseskillicons{trpc}{760}
\showcaseskillicons{truffle}{761}
\showcaseskillicons{tryhackme}{762}
\showcaseskillicons{turborepo}{763}
\showcaseskillicons{turso}{764}
\showcaseskillicons{twig}{765}
\showcaseskillicons{twitch}{766}
\showcaseskillicons{typeorm}{767}
\showcaseskillicons{typescript}{768}
\showcaseskillicons{typst}{769}
\showcaseskillicons{ubuntu}{770}
\showcaseskillicons{uml}{771}
\showcaseskillicons{uniswap}{772}
\showcaseskillicons{unity}{773}
\showcaseskillicons{unitycatalog}{774}
\showcaseskillicons{unocss}{775}
\showcaseskillicons{unrealengine}{776}
\showcaseskillicons{unstructured}{777}
\showcaseskillicons{upwork}{778}
\showcaseskillicons{v}{779}
\showcaseskillicons{vagrant}{780}
\showcaseskillicons{vala}{781}
\showcaseskillicons{vapor}{782}
\showcaseskillicons{vegaspro}{783}
\showcaseskillicons{vercel}{784}
\showcaseskillicons{vim}{785}
\showcaseskillicons{virtualbox}{786}
\showcaseskillicons{visio}{787}
\showcaseskillicons{visualbasic}{788}
\showcaseskillicons{visualstudio}{789}
\showcaseskillicons{vite}{790}
\showcaseskillicons{vitepress}{791}
\showcaseskillicons{vitest}{792}
\showcaseskillicons{vmwareworkstation}{793}
\showcaseskillicons{vscode}{794}
\showcaseskillicons{vscodeinsiders}{795}
\showcaseskillicons{vscodium}{796}
\showcaseskillicons{vuejs}{797}
\showcaseskillicons{vuetify}{798}
\showcaseskillicons{vyper}{799}
\showcaseskillicons{wails}{800}
\showcaseskillicons{wandb}{801}
\showcaseskillicons{warp}{802}
\showcaseskillicons{webassembly}{803}
\showcaseskillicons{webflow}{804}
\showcaseskillicons{webpack}{805}
\showcaseskillicons{websocket}{806}
\showcaseskillicons{webstorm}{807}
\showcaseskillicons{webstudio}{808}
\showcaseskillicons{wezterm}{809}
\showcaseskillicons{windicss}{810}
\showcaseskillicons{windmill}{811}
\showcaseskillicons{windows}{812}
\showcaseskillicons{winedt}{813}
\showcaseskillicons{wireshark}{814}
\showcaseskillicons{word}{815}
\showcaseskillicons{wordpress}{816}
\showcaseskillicons{workers}{817}
\showcaseskillicons{wsl}{818}
\showcaseskillicons{wxt}{819}
\showcaseskillicons{x}{820}
\showcaseskillicons{xcode}{821}
\showcaseskillicons{xd}{822}
\showcaseskillicons{xtable}{823}
\showcaseskillicons{yaml}{824}
\showcaseskillicons{yammer}{825}
\showcaseskillicons{yarn}{826}
\showcaseskillicons{yew}{827}
\showcaseskillicons{yii}{828}
\showcaseskillicons{youtube}{829}
\showcaseskillicons{yui}{830}
\showcaseskillicons{zabbix}{831}
\showcaseskillicons{zed}{832}
\showcaseskillicons{zellij}{833}
\showcaseskillicons{zen}{834}
\showcaseskillicons{zig}{835}
\showcaseskillicons{zudoku}{836}
\showcaseskillicons{zustand}{837}
\end{skilliconcase}


\end{document}
